\documentclass[11pt]{article}
\usepackage[a4paper,margin=1in]{geometry}
\usepackage[T1]{fontenc}
\usepackage{lmodern}
\usepackage{microtype}
\usepackage{parskip}
\usepackage{xcolor}
\usepackage[hidelinks]{hyperref}
\usepackage{titlesec}
\usepackage{enumitem}
\usepackage{amsmath,amssymb}

\definecolor{rcol}{RGB}{40,40,40}
\newenvironment{reviewer}{\color{rcol}\itshape}{\par}
\newcommand{\point}[1]{\vspace{0.7em}\noindent\textbf{#1}\par\vspace{0.15em}}
\newcommand{\reply}{\par\vspace{0.3em}\noindent\textbf{Response.}\ }

\title{Response to reviewers\\[4pt]
\large \emph{Sparse autoencoders reveal organized biological knowledge but minimal regulatory
logic in single-cell foundation models}\\[2pt]
\normalsize Submission ID fbb71648-3418-4139-b1df-d7dd0b5456d6}
\author{Ihor Kendiukhov}
\date{}

\begin{document}
\maketitle

\noindent
I thank the editor and both reviewers for a careful and constructive assessment. The central
criticism---that the manuscript did not separate limitations of the learned representations from
limitations of the input data, the statistical framework, and the evaluation---was well founded,
and addressing it changed the paper substantially. Rather than tempering the wording, I ran the
experiments needed to answer it.

The revision adds twelve new analyses, including two independent external perturbation datasets. Three of them changed headline claims:

\begin{itemize}[leftmargin=1.2em]
\item A capacity-matched comparison against the principal subspace replaces the claim that 99.8\%
of features are ``invisible to SVD''. Truncated SVD needs rank 250--403, not 50, to match the SAE's
own reconstruction quality, and when the leading principal axes are annotated with the identical
protocol they yield at least as many enriched terms per direction as SAE features. The defensible
claim is about parameterization---many sparse, one-sided, individually interpretable directions
rather than a few dense signed ones---not about inaccessible information.

\item A recoverability benchmark establishes what the perturbation data can support at all. The
assay measures 6,546 genes; the median perturbed TRRUST transcription factor has exactly one
curated target among them. Differential expression computed directly on the knockdown---the
strongest perturbation-aware signal these data admit---recovers curated targets for a minority of
factors, and established network-inference methods recover no more. The negative result is now
stated against that ceiling instead of against an implicit optimum of 100\%.


\noindent
\item The regulatory conclusion itself is now qualified. Repeated on perturbation data that measure
the transcriptome, rather than the 6,546-gene panel the atlas was built on, the features recover
targets for 15\% of transcription factors against 1\% for perturbations of non-regulatory genes.
The negative result holds for the primary assay and is too strong as a claim about the models.
\end{itemize}

\noindent
All analyses were also repeated with the cell, rather than the gene position, as the unit of
replication, and the regulatory measurement was replicated in three further cell lines and two
independent external datasets. A summary appears at the end of this letter.

\section*{Editor}

\point{E1. The title should have a clear, precise scientific meaning and should not contain a
colon; it should read as one concise sentence.}
\reply The subtitle has been removed. The title is now
\emph{``Sparse autoencoders reveal organized biological knowledge but minimal regulatory logic in
single-cell foundation models''}, which contains no colon and reads as a single statement.

\point{E2. KEGG copyright permission.}
\reply Permission has been obtained from Kanehisa Laboratories and is being supplied separately to
the editorial office. No KEGG image, pathway map or diagram is reproduced; KEGG is used only as
pathway-to-gene membership lists in the enrichment analysis, and the KEGG references are cited at
every point of use, including in the legends of figures whose content derives in part from KEGG
gene sets.

\point{E3. Code must be deposited in a recognised DOI-assigning repository and linked from Methods
or a dedicated Code Availability section.}
\reply A \textbf{Code availability} section has been added. The complete analysis code, trained SAE
checkpoints and feature catalogs are archived on Zenodo (DOI to be inserted on acceptance) and
mirrored on GitHub.

\point{E4. Key findings insufficiently supported; highlight additional analyses and matched
controls, or tone the conclusions down.}
\reply Both. The additional analyses are listed at the end of this letter; the claims that could
not be supported at their previous strength---the SVD comparison, the cross-layer feature count,
the cross-model comparison, and the scope of the regulatory conclusion---have been restated at the
strength the new evidence supports.

\section*{Reviewer 1}

\point{R1, Major 1. The manuscript interprets low TF--target specificity as evidence of limited
causal regulatory information, but provides no baseline showing how much such information is
recoverable from the same scRNA-seq input.}
\reply This was the most consequential comment and it has been addressed with a new benchmark rather
than with wording.

Every method is now scored in one framework: each emits a ranked list of 100 putative targets per
transcription factor, and each list is tested for enrichment of that factor's curated targets by
hypergeometric test with Benjamini--Hochberg correction across the panel, with the universe set to
the genes the assay actually measures. The comparison includes a perturbation-aware upper bound
(differential expression of the knockdown itself, computed on the same cells), the two standard
network-inference methods the reviewer names in their published formulation (GENIE3 and GRNBoost2,
fitted on the same control cells), Pearson and Spearman co-expression, contextual gene embeddings
taken from the model without any sparse decomposition, the SAE, and size-matched random gene sets.

Under the three reference networks, differential expression of the knockdown recovers targets for
11.1\%, 23.1\% and 44.4\% of factors. GENIE3 reaches 12.5\%, 10.0\% and 0.0\%; GRNBoost2 12.5\%,
10.0\% and 7.1\%; co-expression 25.0\%, 20.0\% and 14.3\%; model embeddings 12.5\%, 10.0\% and
0.0\%. Random gene sets recover nothing, which confirms the test is calibrated.

The benchmark also exposed why the ceiling is low, and this is now reported as a result in its own
right. The assay measures 6,546 genes, and of the 73 perturbed TRRUST transcription factors with
enough cells, only nine have five or more curated targets among them; the median factor has exactly
\textbf{one}. For most of the panel, target-specific recovery is undefined for any method. The
manuscript now states the negative finding against this measured ceiling instead of against an
implicit optimum, and reports the evaluability diagnostics so a reader can see the constraint
directly. Because that constraint is a property of the gene panel rather than of the model, the
measurement was also repeated on two independent datasets with much larger panels (see R1 Major 4).

\point{R1, Major 2. The definition of target genes in the causal-patching analysis is unclear, and
it is not stated whether the evaluation gene set is independent of the genes used for annotation.}
\reply The reviewer identified a real inconsistency, and following it up changed the result.

First, the clarification asked for. The two sections described different things. Features were
selected by an annotation-richness score among features with at least one significant enrichment,
at least ten catalogued top genes and an activation frequency above 0.01; the ``target'' gene set
was then \emph{all 20} of the feature's top activating genes, and ``other'' was every remaining
active position. The evaluation set was therefore not independent of the evidence used to annotate
the feature: ablating a feature was tested on the genes that define it. Both the selection rule and
the gene-set definitions are now stated explicitly in Methods.

Second, the consequence. We repeated the analysis with the evaluation set separated from the
annotation evidence. For each feature the gene positions are partitioned into those both in its
top-20 and in its annotated term; those in the annotated term but \emph{not} in its top-20; and all
others. Each feature also receives its own null, a random gene set matched to its held-out set in
size and mean-activation decile. To rule out the selection rule as a cause, features were drawn from
three arms: the original annotation-richness rule, a uniformly random sample of annotated features,
and a uniformly random sample of all alive features.

On the confounded gene set the effect is large --- median specificity 37.3$\times$, 56.9$\times$
and 66.9$\times$ across the three arms. On the independent evaluation set it is
\textbf{1.04}, \textbf{0.94} and \textbf{0.97}, against matched-random values of 0.89, 0.91 and
0.89. A per-feature Wilcoxon signed-rank test of held-out against matched random gives $p = 0.76$
and $p = 0.34$ in the two unbiased arms. Fourteen of the 119 features tested admit no independent
evaluation at all, because every gene of their annotated term already lies inside their own top-20.

The manuscript now reports this directly: single-feature ablation produces a large, reproducible,
gene-specific effect confined to the genes the feature fires on, and no detectable effect on the
rest of the biological process the feature is annotated with. The previous claim of feature-level
causal specificity for the annotated \emph{process} has been withdrawn, and the Discussion states
the practical consequence for anyone intending to use SAE features as intervention targets.

We note that the analysis also runs on a single stated cell line now (K562). The previously reported
figures came from cells that were all HepG2, for the reason described in the cohort note below.

\point{R1, Major 3. The analyses use only 20 perturbed cells per target for Geneformer and 10 for
scGPT, without power or sample-size sensitivity analyses.}
\reply A sample-size sweep has been added, and it partly vindicates the reviewer's concern while showing
that it does not account for the result.

Cells per target were varied over $n \in \{10, 20, 50, 100\}$ with 20 bootstrap draws each,
restricted to the 14 factors that have at least 100 cells so that sample size is not confounded
with which factors are available. The responding-feature count is flat across the tenfold range
(0.56, 0.65, 0.65, 0.61). Detection of the knocked-down gene itself rises from 5.0\% to 7.1\% and
saturates. Nominal target enrichment improves from 7.7\% at $n{=}10$ to 27.5\% at $n{=}50$ and then
plateaus at 26.9\% for $n{=}100$.

So $n{=}20$ was indeed somewhat under-powered relative to $n{=}50$, and the main analysis now uses
up to 200 cells per target. But the gain saturates well before the effect becomes substantial, and
under correction across the panel no sample size yields a significant transcription factor. The
scGPT arm, which used ten cells per target with position-pooled statistics, has been removed rather
than reported at that power (see the note on scope at the end of this letter).

A positive control was added for the same reason. On identical cells, differential expression
detects the knocked-down gene in 15 of 16 targets at FDR $< 0.05$, with a median Cohen's $d$ of
$-1.19$, while the SAE feature response detects it in 1 of 20. The perturbation is therefore
plainly present in the data and simply is not registered by the feature responses, which separates
a power limitation from a representational one.

\point{R1, Major 4. The regulatory-specificity analysis is limited to K562.}
\reply The measurement has been extended to four cell lines and two independent external datasets.

The CRISPRi resource used here contains four lines, and the analysis now covers all of them, with
perturbed and non-targeting cells always drawn from the same line so that a difference between them
cannot reflect cell-line composition. Pooling the four, \textbf{1 of 66} evaluable transcription
factors shows target enrichment at FDR $< 0.05$: none of 18 in K562, none of 16 in RPE1, one of 16
in Jurkat and none of 16 in HepG2. RPE1 is a non-transformed retinal pigment epithelial line and
HepG2 a hepatocyte line, so the finding is not confined to the leukaemic backgrounds in which these
models are usually probed. The design effect for the unit-of-replication correction is 1.61--1.73
in every line, and in no line do the transcription factors separate from the
non-transcription-factor controls.

Because the gene panel rather than the cell type is the binding constraint on evaluability, the
measurement was also repeated on two perturbation datasets from other groups with other protocols
and much larger panels: ECCITE-seq of THP-1 monocytes~\cite{papalexi2021eccite} (18,649 genes) and
Perturb-seq of K562~\cite{norman2019exploring} (33,694 genes). These raise the evaluability ceiling
substantially --- in the THP-1 data the median perturbed factor has 19.5 measured TRRUST targets
against one in the primary resource --- and therefore test the finding where detection is actually
possible. And they change the answer.

On the THP-1 data the SAE recovers targets for 20--22\% of factors under DoRothEA, and on the
Perturb-seq data 23--25\%, where on the primary panel it recovered none. Whether that recovery is
factor-specific depends on the reference network, and we added the control that decides it: ten
perturbations of genes absent from every reference network, run through the identical pipeline,
with their predicted sets tested against the \emph{same} regulons as the transcription factors.
Against the high-confidence DoRothEA subset a factor's own prediction is enriched for its own
regulon in 3 of 13 cases (23\%); of the nine control genes that produced a prediction, one enriches
for at least one regulon, giving 1 significant test in 117 factor-by-regulon comparisons. The
direction is what specificity predicts, but the panel cannot establish it: Fisher's exact test on
3 of 13 against 1 of 9 gives $p = 0.62$. Against the union of all confidence levels the comparison
reverses --- 4 of 16 factors against 8 of 9 control genes ($p = 0.004$, controls higher) --- because
those regulons have a median of 402 measured genes and a broad predicted set enriches for them
regardless of what was perturbed. That tier does not measure regulatory recovery and we do not
report it as such.

This qualifies our original claim in both directions, and we have been careful not to overcorrect.
The absence of detectable specificity in the primary CRISPRi panel is substantially a property of
that assay rather than of the models: with a median of one measurable target per factor, no method
can demonstrate specificity there. On data that can detect it, the flat null does not reproduce ---
but the one comparison capable of demonstrating a factor-specific signal rests on thirteen factors
in a single screen and does not reach significance. The manuscript now states exactly this, in the
abstract, the Results and the Discussion, and identifies panel size as the binding limitation. We are grateful for the push that led
us to run these datasets; the paper is more accurate for it.

\point{R1, Minor 1. The random seed used for SAE training is not reported, and run-to-run
variability is not assessed.}
\reply The seed is now reported, and the stability analysis the reviewer asks for produced a result that
we think belongs in the paper rather than only in this letter.

The atlas dictionaries were trained with random seed 42, which also fixes the training subsample.
This is now stated in Methods. Some care was needed in reporting it: the previous training script
seeded NumPy before selecting the training subsample but did not seed the PyTorch generator that
initialises the weights, so initialisation and batch order were not in fact controlled by that seed.
The released code seeds both explicitly and records the seed in every output file.

For the stability analysis we retrained at layers 0 and 11 under six runs each: five differing only
in weight initialisation and batch order with the training subsample held fixed, and one with a
different 1M-position subsample.

Aggregate properties are highly reproducible. Variance explained is $0.8454 \pm 0.0003$ at layer~0
and $0.8200 \pm 0.0003$ at layer~11; the annotation rate is $0.935 \pm 0.009$ and $0.933 \pm 0.010$;
the module count is $7.2 \pm 0.4$ and $8.3 \pm 1.0$; dead-feature counts are $2 \pm 1$ and
$32 \pm 4$. Every quantity the manuscript reports at the level of the atlas is stable to retraining.

Individual dictionary atoms are not. The mean best decoder cosine between runs is 0.282 at layer~0
and 0.294 at layer~11, and fewer than 1\% of atoms have a counterpart above 0.9 in another run. To
make those numbers interpretable we computed the reference empirically: the best match a random
unit direction finds among 4,608 directions in 1,152 dimensions is $0.113 \pm 0.008$, and no random
direction reaches 0.9. Cross-seed agreement is therefore well above chance and far below
reproducibility.

We report this as a limitation in the manuscript rather than only here, because it bears directly on
how the released atlas should be used: aggregate structure is a property of the model, while the
identity of any individual catalogued feature is partly a property of the training run. It is also
consistent with the gene-level program analysis added for R2 Major 1, in which 82,525 atoms resolve
into 18,711 distinct programs --- the programs appear to be the level at which the representation is
stable, not the atoms.

\point{R1, Minor 2. The enrichment-term labels in Figure 6 are truncated and several axis labels
overlap.}
\reply Confirmed and fixed. In the previous version the enrichment labels in the causal-patching
figure were hard-truncated at 20 characters, which cut four GO accessions mid-identifier and
collapsed two label pairs into identical strings, and the neighbouring panel's tick labels
overprinted. The figures have been regenerated: labels are wrapped rather than truncated, tick
density is set per panel width, and no text in any figure is now clipped or overlapping. All
figures are drawn directly from the result files by a single script included in the code archive,
so a figure cannot drift from the numbers it depicts.

\point{R1, Minor 3. Many internal cross-references remain unresolved and are displayed as
``Table ??''.}
\reply Fixed at the root cause. The previous submission split the main text and the supplementary
information into two LaTeX documents that imported each other's labels, which resolves correctly
only if the two files are compiled in a specific alternating order; a single-pass build renders
every such reference as ``??''. The main text no longer depends on that mechanism: supplementary
items are referenced by explicit number, and a check script included with the source verifies every
such reference against the supplementary document and fails the build if any is missing or
misnumbered. The compiled PDFs contain no unresolved references.

\section*{Reviewer 2}

\point{R2, Major 1. Summing 82,525 features across layers and inferring more than 70-fold
compression into a 1,152-dimensional space is not conceptually justified: each layer has its own
space and its own independently trained dictionary.}
\reply The reviewer is correct, and the claim has been removed. Dictionary atoms trained at
different layers occupy different vector spaces and cannot be summed into a count of concepts
coexisting in one representation.

To replace it with something measurable, features are now compared in the space they do share---the
gene space---through their top-20 gene signatures. Clustering all 82,525 Geneformer features by
gene-set similarity yields \textbf{18,711 distinct gene-level programs} (stable across Leiden
resolutions 0.5--2.0), that is 4.4 atoms per program, with 97\% of multi-atom programs recurring in
two or more layers. Within a single layer, which is the only place a features-per-space ratio is
defined, the dictionary is \textbf{4$\times$ overcomplete} (4,608 atoms in 1,152 dimensions) and
resolves roughly \textbf{1,250 distinct programs}. The manuscript now reports these quantities and
the packing geometry that bounds them (mean decoder coherence 0.036 against a Welch bound of
0.0255), and states the aggregate count explicitly as an aggregate over 18 separately trained
dictionaries.

\point{R2, Major 2. The SAE--SVD comparison is not capacity-matched: 4,608 features against 50
singular directions, with an arbitrary cosine threshold. The comparison should be repeated under
capacity- or reconstruction-matched conditions, ideally by quantifying the cumulative projection of
each SAE direction onto the top-$k$ SVD subspace across a range of $k$.}
\reply Done exactly as suggested, and the result changed the claim.

For every layer we computed the exact eigendecomposition of the activation covariance and, for each
unit-normalized decoder direction, the cumulative projection $\rho_k$ onto the rank-$k$ principal
subspace for all $k$ up to the full dimension, against the random-direction reference
$\mathbb{E}[\rho_k] = k/d$. Three findings follow.

\emph{Reconstruction-matched capacity.} Truncated SVD requires rank \textbf{250--403} to match the
SAE's variance explained (for example, rank 308 at layer 11 against the SAE's 0.823). The previous
comparison was five- to eightfold under-capacity.

\emph{Projection.} The median SAE direction places 50\% of its norm within the leading 152--247
components and 90\% within the leading 664--764, compared with 576 and 1,037 for a random
direction. SAE directions are two- to fourfold more concentrated in the leading principal subspace
than chance---they are not hidden from it.

\emph{Alignment is not a capacity artefact.} The fraction of features aligned with any single axis
above a cosine of 0.5 saturates by $k \approx 20$--50 and does not rise thereafter: even against
the \emph{complete} 1,152-dimensional basis it is 0.09--0.98\% depending on layer. Low single-axis
alignment therefore reflects that the information is distributed over many axes, not that too few
axes were offered.

We also annotated the leading 50 principal axes with the identical top-20-gene enrichment protocol
used for SAE features (both polarities, since a principal axis is signed whereas a TopK feature is
one-sided). Per direction they are at least as annotatable as SAE features: 98\% annotated with
159--179 enriched terms per direction, against 90--93\% and 64--71 terms for randomly drawn SAE
features. The claim that biological signal is exclusive to non-aligned features does not survive a
matched comparison and has been removed.

The manuscript now states what the matched comparison does support: the dictionary provides a
different parameterization of largely the same subspace---many sparse, one-sided, individually
interpretable directions in place of a few dense, signed, polysemantic ones---which is what makes
per-feature causal intervention and per-feature annotation possible.

\point{R2, Major 3. Potentially serious pseudoreplication: roughly 20 cells per target, but tests
against 100K control positions, treating gene positions from the same cell as independent
replicates.}
\reply The reviewer is right that the unit of replication was wrong, and the analysis has been redone with
the cell as the unit throughout. We also measured how large the effect of the error was, rather than
assuming it.

The intraclass correlation of feature activation across gene positions within a cell is
$2.98 \times 10^{-4}$. At the observed 2,020 positions per cell this gives a design effect of
\textbf{1.61}, so the effective sample size is about 62\% of the nominal position count rather than
smaller by three orders of magnitude. TopK sparsity leaves activations at different positions of one
cell close to independent, which is why the inflation is bounded. The same quantity is 1.66--1.73 in
the other three cell lines, so this is not particular to one dataset.

The correction has been applied regardless. Significance is now assessed by a Mann--Whitney test
over cells with Benjamini--Hochberg correction across features, and the effect-size definition is
held fixed at the control position-level scale so that the cell-level and position-level results
differ \emph{only} in what is treated as an independent observation. On that basis the two are
close (a mean of 1.90 responding features per factor over cells against 1.60 over positions),
consistent with the measured design effect.

A larger statistical problem surfaced while making this correction. The previously reported response
counts were thresholded on effect size without any correction for testing 4,608 features.
Reproducing that criterion gives a mean of 3.20 responding features per factor with 15 of 20 factors
showing at least one; applying Benjamini--Hochberg across features collapses it to 2 of 20. The
manuscript now reports both, and a negative control has been added: ten perturbations of genes that
are not transcription factors yield \emph{more} responding features (2.30) than the twenty
transcription factors (1.90), so what the features register is not specific to
transcription-factor perturbation at all.

\point{R2, Major 4. Substantial confounding in the cross-model comparison: Geneformer activations
from K562, scGPT activations from Tabula Sapiens.}
\reply The confound was real, and removing it reversed the conclusion.

We trained Geneformer dictionaries on the same 3,000 Tabula Sapiens cells used for the scGPT atlas,
with identical architecture and hyperparameters, and compared the two models at matched relative
depth (layer index divided by depth minus one). Because the published scGPT dictionaries were
trained on all available positions rather than a 1M subsample, we also retrained them under the
Geneformer protocol as a control; the conclusion is the same either way.

Variance explained on held-out positions:

\begin{center}
\begin{tabular}{lccc}
\hline
relative depth & Geneformer / TS & scGPT / TS & scGPT / TS (matched protocol) \\
\hline
0.00 & \textbf{0.7731} & 0.7258 & 0.7676 \\
0.28 & 0.7749 & 0.7729 & 0.7827 \\
0.65 & 0.7596 & \textbf{0.8515} & 0.8171 \\
1.00 & 0.7241 & \textbf{0.8862} & 0.8552 \\
\hline
\end{tabular}
\end{center}

The manuscript had reported a uniform scGPT advantage in reconstruction (means of 90.2\% against
81.7\%) and read it as architectural. On matched inputs there is no uniform advantage: Geneformer
reconstructs better at the input end, the two are indistinguishable at about a quarter depth, and
scGPT is ahead through the second half. The difference is a depth-dependent crossover, which is not
visible unless the inputs are matched.

The size of the confound is easy to state. The same Geneformer dictionaries reconstruct their own
CRISPRi-like data 5--8 percentage points better than Tabula Sapiens at every depth (0.8444 against
0.7731 at layer~0, for example), which is the same order as the cross-model gap that had been
attributed to architecture.

The cross-model claims in the manuscript have been rewritten accordingly, and the comparison table
now reports matched-input values at matched relative depth. Decoder coherence is 0.030--0.036 for
Geneformer and 0.039--0.047 for scGPT against Welch bounds of 0.0255 and 0.0362, so both
dictionaries sit close to the packing limit for their dimensionality and neither is exploiting
unusual capacity.

\point{R2, Major 5. The conclusions extend beyond the evidence; claims in the title, abstract and
discussion should be restricted to the cellular context, perturbation dataset, SAE design and
regulatory reference evaluated here, and the analysis extended to more cell types.}
\reply We agree that a negative result obtained under one analytical framework cannot by itself establish
a property of a model class, and the revision addresses this by widening the evidence rather than
by qualifying the sentences.

The regulatory measurement now spans four cell lines --- two of them non-leukaemic --- and two
perturbation datasets generated by other groups with other protocols and much larger gene panels,
with a perturbation-aware ceiling computed on every one so that the recovery rate is interpretable
rather than merely low. The scope that genuinely constrains the claim is stated where the claim is
made, in the abstract and the Discussion: what was tested, on which data, against which reference
networks.

On the specific question of whether the models lack regulatory logic or our pipeline fails to
detect it, we have replaced argument with measurement. On identical cells, differential expression
detects the knocked-down gene in 15 of 16 targets while the SAE feature response detects it in 1 of
20. The signal is present and is not registered. That locates the limitation precisely --- the
features track global cell-state shifts rather than gene-level consequences of the intervention ---
and the Discussion states the conclusion in those terms.

We would add one point of principle. No finite set of probes can establish that regulatory
information is absent from a model, because the proposition that it resides in some representation
not yet examined is not refutable by examination. What can be done is to enumerate the
representations that have been examined and report what was found in each. The Discussion now does
this explicitly, across attention-derived edge scores, causal circuit tracing, exhaustive circuit
mapping, causal intervention, adversarial validation of in-silico perturbation profiles,
representation geometry, and the sparse decomposition reported here.

\point{R2, Minor 1. Incomplete cross-references.}
\reply Addressed at the root cause; see R1 Minor 3 above.

\point{R2, Minor 2. Several references appear inaccurate or incompletely formatted, including
references 20 and 23.}
\reply The reviewer is right about both, and about others. The full bibliography was re-verified
entry by entry against publisher records.

References 20 and 23 were the most serious: both carried a placeholder author field
(``and others'') and an editorial note that was being typeset into the reference list itself, and
both had the wrong venue and an inexact title. They are now
Wang \emph{et al.}, \emph{Nature Communications} \textbf{17}, 730 (2025),
doi:10.1038/s41467-025-67481-2, and
Wang \emph{et al.}, \emph{Bioinformatics} \textbf{41}(7), btaf363 (2025),
doi:10.1093/bioinformatics/btaf363, with complete author lists and the notes removed.

The audit also corrected: the TRRUST page range (D199--D205 $\rightarrow$ D380--D386) and an author
name; the STRING page range (D599--D610 $\rightarrow$ D638--D646) and a truncated author list; an
author-name error in the Geneformer reference (Manber $\rightarrow$ Mantineo); an author-name error
in the Replogle reference (Alexandra $\rightarrow$ Angela Pogson); the venue for the sparse-probing
reference (arXiv $\rightarrow$ ICLR 2024) with the published author order; and author-name errors in
two further entries. Web-only sources are now proper \texttt{@misc} entries with URLs and access
years, and a DOI has been added to every one of the 19 entries that has one.

\section*{A correction to the description of the data}

While repeating the causal-ablation analysis we found an error in how the previous version
described its own cohort, and we want to report it plainly.

The manuscript stated that the atlas was built from ``2,000 K562 control cells''. It was not. The
extraction selects control cells by perturbation label alone (\texttt{gene == "non-targeting"})
across the whole processed Replogle matrix, which spans four cell lines, and applies no cell-line
filter. The cohort is 591 Jurkat, 581 K562, 567 RPE1 and 261 HepG2 cells. We verified this rather
than inferring it from the code: re-running the selection under its recorded seed reproduces the
per-cell gene counts stored at extraction time exactly.

Two consequences follow. The causal-ablation analysis took the first 200 cells of that cohort after
sorting by row index, and because the matrix is ordered by cell line, those 200 cells were
\emph{all HepG2}. And the perturbation analysis compared K562 perturbed cells against a control
baseline drawn from all four lines, so a cell-line difference could contribute to an apparent
perturbation response.

Nothing about the atlas itself needed recomputing --- the dictionaries, annotations, modules,
cross-layer structure and geometry are all properties of the cells that were actually processed,
and they stand. What was wrong was the description. The manuscript now states the cohort
composition explicitly, reports it in Supplementary Table~S1, and refers to the multi-line control
atlas rather than attributing it to any single line. Every analysis that compares perturbed against
control cells now draws both from the same line, and causal ablation is run within one stated line.

We are grateful that the reviewers' scrutiny of the causal-ablation design led us to this.

\section*{A note on scope}

Two analyses present in the previous version have been removed rather than revised, and we want to
be explicit about it.

The scGPT causal-patching result was obtained by feeding the model a uniform value bin for every
gene instead of the per-cell binned expression values, because the value-embedding inputs had not
been retained alongside the hidden states. The scGPT perturbation-specificity result used ten cells
per target with position-pooled statistics. Now that every other measurement in the paper is
assessed over cells, uses the actual model inputs, and is corrected for multiple testing, reporting
those two at their original standard would be inconsistent with the rest; and bringing them up to
that standard would require re-extracting the scGPT activations, which the present software
environment does not support alongside the rest of the pipeline.

We have therefore confined the regulatory analysis to Geneformer and say so in the Limitations. The
scGPT \emph{atlas} is unaffected and is reported in full: dictionary statistics, ontology
annotation, co-activation modules, cross-layer connectivity, cell-type enrichment, and the batch
and assay analysis. What changes is that no regulatory claim is now made on the basis of the scGPT
arm.

\section*{Summary of new analyses}

\begin{center}
\begin{tabular}{p{4.4cm}p{10.4cm}}
\hline
\textbf{Analysis} & \textbf{What it establishes} \\
\hline
Cohort composition &
The control cohort spans all four cell lines in the resource, not one. Reported explicitly and
verified against the extraction metadata. \\
Capacity-matched principal-subspace comparison &
Reconstruction-matched SVD rank, cumulative projection curves against a random-direction
reference, alignment as a joint function of subspace size and threshold, and annotation yield per
direction. \\
Distinct-concept counting &
Features compared in the shared gene space; distinct gene-level programs per layer and in
aggregate; packing geometry against the Welch bound. \\
Recoverability benchmark &
Every method scored in one framework: differential expression of the knockdown, GENIE3, GRNBoost2,
co-expression, model embeddings, the SAE, and random controls. \\
Assay evaluability &
How many perturbed transcription factors have enough curated targets among the measured genes to
be evaluable at all. \\
Cell-level statistics &
Intraclass correlation and design effect measured rather than assumed; significance reassessed
with the cell as the unit of replication at a fixed effect-size definition. \\
Non-transcription-factor controls &
Perturbations of genes that are not transcription factors, as a negative control for the
specificity test. \\
Sample-size sweep &
Specificity and a positive control as a function of perturbed cells per target, with bootstrap
intervals. \\
Replication across cell lines &
The regulatory measurement repeated in RPE1, Jurkat and HepG2, with perturbed and control cells
drawn from the same line. \\
Independent external datasets &
The same measurement on ECCITE-seq of THP-1 monocytes and on a near-transcriptome-wide Perturb-seq
of K562, which have far larger measured gene panels. \\
Causal ablation with independent evaluation genes &
Held-out term genes, per-feature expression-matched random controls, and unbiased feature
selection alongside the original selection rule. \\
Run-to-run stability &
Five independent training seeds per layer plus a different training subsample; dictionary
agreement and the spread of every reported quantity. \\
Matched cross-model design &
Geneformer dictionaries trained on the same cells as the scGPT atlas, compared at matched relative
depth. \\
\hline
\end{tabular}
\end{center}

\end{document}
